%!TEX root = ../main_thesis.tex

\chapter{Il linguaggio ROOC}
\label{cap:linguaggio}

\todo[inline]{Capitolo di riferimento del linguaggio, dal punto di vista dell'utente: cosa si può scrivere e qual è il significato dei costrutti. L'implementazione è rimandata ai capitoli successivi.}

\section{Principi di design}
\label{sec:design-linguaggio}

\todo[inline]{Illustrare gli obiettivi di progetto: modello formale, leggibilità, vicinanza alla notazione matematica, separazione tra modello e dati.}

\section{Struttura di un modello}
\label{sec:struttura-modello}

Un modello ROOC è composto dalla funzione obiettivo (\texttt{min}/\texttt{max}),
dai vincoli (\texttt{s.t.}) e da due blocchi opzionali: \texttt{where} per la
definizione di costanti e \texttt{define} per la dichiarazione del dominio delle
variabili.

\todo[inline]{Mostrare la struttura generale con uno snippet commentato. Spiegare obiettivo, vincoli con iterazione (\texttt{for}), blocco \texttt{where} e blocco \texttt{define}.}

\section{Esempio guida}
\label{sec:esempio-guida}

\todo[inline]{Presentare per esteso uno o due esempi (knapsack binario e dominating set su grafo) che verranno richiamati nei capitoli successivi per illustrare le varie fasi di compilazione.}

\section{Primitive e tipi di dato}
\label{sec:primitive}

Il linguaggio offre un insieme ricco di tipi: numeri, stringhe, booleani,
\emph{grafi}, tuple e iterabili. Si descrive la semantica di ciascuno.

\todo[inline]{Descrivere le primitive (\texttt{Number}, \texttt{Integer}, \texttt{String}, \texttt{Boolean}, \texttt{Graph}, \texttt{Tuple}, \texttt{Iterable}) e la loro rappresentazione a livello di tipo. Riferimento: \texttt{src/primitives/}.}

\section{Iteratori, \emph{block functions} e destructuring}
\label{sec:iteratori}

\todo[inline]{Spiegare gli iteratori (\texttt{for ... in ...}), le \emph{block scoped functions} (\texttt{sum}, \texttt{prod}, \texttt{min}, \texttt{max}, \texttt{avg}) e il destructuring di tuple e archi di grafo.}

\section{Operatori e definizione delle variabili}
\label{sec:operatori-variabili}

\todo[inline]{Trattare gli operatori e l'\emph{operator overloading} sulle primitive, e la dichiarazione del dominio delle variabili con i relativi \emph{bounds} (Boolean, Real, NonNegativeReal, IntegerRange).}

\section{Libreria standard e funzioni definite dall'utente}
\label{sec:funzioni}

\todo[inline]{Elencare le funzioni built-in (su array, grafi, range) e descrivere il meccanismo di estensione con funzioni definite dall'utente (anche via JavaScript in ambiente WASM). Riferimento: \texttt{src/runtime\_builtin/}.}
